\chapter[Introduction]{When are DAOs Dead? An Empirical Model for Determining the Activity of a Decentralized Autonomous Organization}

\noindent In November 2021, a group of cryptocurrency enthusiasts raised \$47 million in less than a week to purchase an original copy of the United States Constitution. Despite failing to win the auction, ConstitutionDAO demonstrated the remarkable potential of blockchain-based organizations to rapidly mobilize resources and coordinate action. Yet today, while the organization's token maintains a market value exceeding \$250 million, ConstitutionDAO itself is effectively defunct -- highlighting a fundamental challenge in evaluating these novel organizational forms. When exactly does a Decentralized Autonomous Organization (DAO) transition from active to inactive? This question has profound implications for investment decisions and the future of decentralized governance.

A Decentralized Autonomous Organization represents a new form of organizational structure. It operates through rules encoded by computer programs, transparent and verifiable on a blockchain. Unlike traditional organizations with hierarchical management structures, DAOs implement governance mechanisms through smart contracts which are self-executing programs that automatically enforce organizational rules and decisions (\cite{mik2017smart}). Member participation occurs through token-based voting systems, where governance tokens grant holders the right to propose and vote on organizational decisions. This technological infrastructure enables DAOs to operate using a flat hierarchy, relying on transparent, programmable rules that all participants can verify (\cite{el2020overview}, \cite{witt2024blockchain}).

Traditional metrics for organizational activity -- revenue, employee count, stock performance -- prove inadequate for evaluating these blockchain-native organizations, as evidenced by the passing of Wyoming's Decentralized Unincorporated Nonprofit Association Act (\cite{kerr2024legal}). Since \textcite{buterin2013ethereum} introduced the concept in the Ethereum whitepaper, DAOs have evolved into a significant force in the digital economy. These organizations collectively manage treasuries worth \$37.8 billion according to DeepDAO, a DAO aggregation platform. Leading examples demonstrate their impact: Uniswap, a protocol to swap two on-chain assets, oversees \$3.2 billion in assets, while Optimism, a layer two implementation, controls over \$2 billion. Beyond pure financial scale, DAOs are also tackling real-world challenges -- GitcoinDAO has distributed more than \$50 million to open-source projects, while KlimaDAO has facilitated the retirement of millions of tons of carbon credits.

A defining characteristic of DAOs is their unprecedented accessibility and ease of formation. Unlike traditional organizations that require substantial legal infrastructure, regulatory compliance, and formal registration processes, DAOs can be created by any individual or group with basic technical knowledge and minimal capital investment. This low barrier to entry has led to a proliferation of DAOs across the blockchain ecosystem. The technical infrastructure provided by platforms like Aragon and DAOhaus enables users to deploy standardized organizational templates within minutes, requiring only the specification of basic governance parameters and initial token distribution.

However, the characteristics which enable a DAOs rapid formation and flexible operation create significant methodological challenges for researchers. Unlike traditional organizations, which must formally dissolve to become inactive, DAOs can enter a state of operational abandonment while maintaining a technical existence on the blockchain (\cite{wright2015decentralized}). This distinctive property presents a fundamental research challenge: How can scholars effectively distinguish between active and inactive DAOs when conventional organizational metrics prove insufficient? The challenge becomes particularly acute given that 72\% of DAOs become operationally dormant shortly after creation (\cite{alphaquest2024dead}), yet persist indefinitely due to the immutable nature of blockchain infrastructure.

Current research approaches demonstrate the scope of this methodological challenge. Studies such as \textcite{monteiro2024exploring}'s analysis of 179 DAOs rely on carefully curated samples of known active organizations, introducing potential survivorship bias that limits our ability to draw broader conclusions about organizational dynamics. Other attempts to classify DAO activity levels, such as \textcite{bellavitis2023rise}, define active organizations through governance participation metrics. However, this one dimensional approach may fail to capture the different operational states that DAOs may occupy. For instance, mature DAOs may achieve a state where core functionalities fully satisfy organizational objectives, resulting in reduced governance activity despite continued operational effectiveness. Furthermore, the rapid evolution of the DAO landscape makes it increasingly difficult to establish representative sample parameters -- the population characteristics shift continuously as new organizations form and others cease operations.

Existing measurement frameworks show significant limitations in addressing these challenges. While platforms like DeepDAO offer comprehensive metrics including governance participation rates, they struggle to establish clear thresholds for determining organizational activity levels. Similarly, CoinGecko's classification methodology lacks transparency, impeding scholarly validation and replication of their findings (\cite{Lee_2024}). These limitations create substantial barriers to conducting large-scale, representative studies of DAO organizational dynamics.

To address these methodological gaps, our research pursues two primary objectives: (1) to evaluate the reliability of simplified measurement approaches based on individual organizational metrics in predicting DAO activity levels, and (2) to develop and validate a machine learning model that synthesizes financial, social, and technical indicators to accurately classify organizational activity states. Through these objectives, we seek to determine whether multidimensional analysis offers meaningful advantages over simplified approaches in assessing DAO operational status.

The practical implications of this research extend across multiple stakeholder groups within the blockchain ecosystem. For private and institutional investors, our methodology provides empirically validated tools for evaluating project viability before capital deployment. DAO administrators gain access to standardized benchmarking metrics for assessing organizational performance against industry norms. Researchers benefit from a systematic methodology for sample selection, enabling more reliable comparative studies of decentralized organizations. Last but not least, regulatory bodies and policymakers gain access to more precise analytical tools for understanding DAO ecosystem dynamics, supporting evidence-based policy development.


The remainder of this thesis systematically addresses these research objectives through five chapters. Chapter 2 establishes the theoretical foundations of decentralized organizational forms and presents our research hypotheses, examining how monetary, social, and technical characteristics inform our understanding of DAO activity patterns. Chapter 3 details our methodological approach, describing the development of a comprehensive dataset encompassing 1,332 DAOs and outlining both our heuristic benchmarks and machine learning model specifications. Chapter 4 presents our empirical analysis, comparing simplified heuristic approaches against our machine learning model, with results demonstrating 75\% accuracy in classifying DAO activity states. Chapter 5 discusses theoretical and practical implications, with particular emphasis on methodological applications for researchers studying decentralized organizations. Finally, Chapter 6 synthesizes key findings and outlines promising directions for future research in organizational activity measurement within decentralized contexts, while acknowledging the limitations of our analytical approach.